\section{Variational Derivation}
\label{app:variational}
\label{sec:variational}

We derive the governing equations from Hamilton's extended principle, following Bedford \& Drumheller~\cite{drumheller1980thermomechanical} and Drumheller~\cite{drumheller2000}, adapted to a partially saturated four-phase system (solid, aquous, oelic, gas) and arbitrary miscible components $\alpha$ within each phase and with porosity-dependent solid free energies.

Each component has volume fraction~$\phi_\xi$, intrinsic density~$\bar\rho_\xi$, bulk density $\rho_\xi = \phi_\xi\bar\rho_\xi$, velocity~$\mbf{v}_\xi$, deformation gradient~$\mbf{F}_\xi$, and Jacobian $J_\xi = \det(\mbf{F}_\xi)$.  The solid phases share a common trajectory ($\mbf{v}_A = \mbf{v}_B = \mbf{v}_s$, $\mbf{F}_A = \mbf{F}_B = \mbf{F}_s$, $J_A = J_B = J_s$).  The total porosity is $\phi = \phi_f + \phi_g$, the fluid saturation is $S_f = \phi_f/\phi$, and the mass added per unit reference volume to component~$\xi$ by reactions is~$c_\xi$, with rate per unit current volume $c^+_\xi = \dot{c}_\xi / J_\xi$.

The free energies of the solid phases are written as $\psi_A(\mbf{F}_s, \phi)$ and $\psi_B(\mbf{F}_s, \phi)$, retaining the full deformation gradient and the current porosity as independent arguments.  The multiplicative decomposition $\mbf{F}_s = \mbf{A}_s \bar{\mbf F}_s$ is \emph{not} imposed at this stage; it enters later in the Coleman--Noll procedure (Appendix~\ref{sec:coleman}) when a specific constitutive form is adopted.

\subsection{Variational principle and its terms}

The extended Hamilton's principle is
\begin{equation*}
  \int_{t_1}^{t_2}\left[\delta(T - U) + \delta W + \sum_k \delta C_k\right] dt = 0,
\end{equation*}
with kinetic energy, internal energy, virtual work, and three constraints:
\begin{align*}
  T &= \sum_\xi \int_{\mc{B}} \tfrac{1}{2}\rho_\xi\,\mbf{v}_\xi\cdot\mbf{v}_\xi \,\dv,  \\
  U &= \sum_\alpha \left[\int_{\mc{B}_{s0}} \rho_{s0}\,\psi^\alpha_s(\mbf{F}_s, \bar\rho^\alpha_s, \eta^\alpha_s,  \theta^\alpha_s, \phi) \,  dV_s \right. \\
    &\qquad + \int_{\mc{B}_{a0}} \rho_{a0}\,\psi_a^\alpha(\bar\rho_a^\beta, \eta_a^\beta, \theta^\beta_a)\, dV_a + \int_{\mc{B}_{o0}} \rho_{o0}\,\psi^\alpha_o(\bar\rho_o^\beta, \eta_o^\beta, \theta^\beta_o)\,dV_o +  \int_{\mc{B}_{g0}} \rho_{g0}\,\psi^\alpha_g(\bar\rho^\alpha_g, \eta_g^\beta, \theta^\beta_g)\,dV_g \\ 
    &\left.\qquad + \int_{\mc{B}} \gamma_{oa}(S_a) \,\dv + \int_{\mc{B}} \gamma_{go}(S_g)\,\dv \right],  \\
  \delta W &= \sum_\xi \int_{\mc{B}} \left[(\rho_\xi\mbf{g} + \mbf{b}_\xi)\cdot\delta\mbf{x}_\xi + \sum_\alpha \mc{L}^\alpha_\xi\,\frac{\delta c^\alpha_\xi}{J_\xi}\right] \,\dv,  \\
  C_1 &= \int_{\mc{B}} \lambda\left(\sum_\xi\phi_\xi - 1\right) \,\dv,  \\
  C_2 &= \sum_\xi \sum_\alpha \int_{\mc{B}_{\xi 0}} \mu_\xi^\alpha \left( J_\xi \eta_\xi^\alpha - \frac{\phi_{\xi 0} \bar\rho_{\xi 0}^\alpha + c_\xi^\alpha}{\phi_\xi \bar\rho_\xi^\alpha} \right) dV_\xi \\ 
  C_3 &= \sum_\xi \int_{\mc{B}_{\xi 0}} \pi_\xi \left( \sum_\alpha \eta_\xi^\alpha - 1 \right) dV_\xi,  \\
  C_4 &= \int_{\mc{B}} \tau\sum_\alpha\sum_\xi \frac{\dot{c}^\alpha_\xi}{J_\xi} \,\dv.
\end{align*}
The Helmholtz free energies are $\psi^\alpha_s$, $\psi_a^\alpha$, $\psi_o^\alpha$, and $\psi_g^\alpha$.  The interaction forces satisfy $\sum_\xi \mbf{b}_\xi = \mbf{0}$, and $\mc{L}^\alpha_\xi$ is a generalized force conjugate to mass transfer (identified later via Coleman--Noll).  The Lagrange multipliers~$\lambda$, $\mu^\alpha_\xi$, $\pi_\xi$, and~$\tau$ enforce the volume fraction constraint, component mass conservation, component partitioning, and mixture mass balance, respectively.

  We treat $\mbf{x}_\xi$, $\eta^\alpha_\xi$, $\phi_\xi$, $\bar\rho^\alpha_\xi$, and $c_\xi^\alpha$ as independent fields and compute the variations in what follows. Most of the required variations are detailed in Bedford \& Drumheller~\cite{drumheller1980thermomechanical} and will be omitted here; however, there are some new terms in this work that are derived below.

\subsection{Variation of the internal energy}

The solid contribution is
\begin{align}
  \delta U_s &= \int_{\mc{B}_{s0}} \sum_{\xi=A,B} \rho_{\xi 0} \left[\left.\p{\psi_\xi}{F_{kl}}\right\vert_\phi \delta F_{kl} + \left.\p{\psi_\xi}{\phi}\right\vert_{\mbf{F}_s} \delta\phi \right] dV_s.
  \label{eq:deltaU_solid}
\end{align}
The variation $\delta F_{kl} = \partial(\delta x_{s,k})/\partial X_{s,l}$ produces the standard stress divergence upon integration by parts.  Transforming to the current configuration gives
\begin{equation}
  \int_{\mc{B}_{0}} \rho_{\xi 0}\,\left.\p{\psi_\xi}{F_{kl}}\right\vert_\phi \delta F_{kl}\, dV_s = -\int_{\mc{B}} \nabla_{\mbf{x}} \cdot (\bs\sigma'_\xi)^T \cdot \delta\mbf{x}_s \,\dv + \text{boundary terms},
  \label{eq:stress_variation}
\end{equation}
where the effective Cauchy stress is defined through
\begin{equation}
  \bs\sigma'_\xi = \frac{1}{J_s} \rho_{\xi 0}\,\p{\psi_\xi}{\mbf{F}_s}\bigg\vert_\phi\, \mbf{F}_s^T.
  \label{eq:var_cauchy_def}
\end{equation}
The notation $|_\phi$ emphasizes that this derivative is taken at fixed porosity.

\subsection{Porosity variation via the inter-component identity}

The porosity~$\phi = 1 - \phi_s$ is a spatial field naturally associated with the fluid and gas components, yet it appears as an argument of the solid free energies $\psi^\alpha_s(\mbf{F}_s, \phi, \ldots)$ in the integral over the solid reference configuration.  Since this integral is taken at fixed~$\mbf{X}_s$, we require the variation of~$\phi$ at fixed~$\mbf{X}_s$.  Using the inter-component variation identity (equation~(1.26) of Bedford \& Drumheller~\cite{bedford1983theories})
\begin{equation}
  \left.\delta\Gamma_{(\zeta)}\right|_{\mbf{X}_{(\xi)}} = \delta\Gamma_{(\zeta)} - \nabla_{\mbf{x}}\Gamma_{(\zeta)} \cdot \left(\delta\mbf{x}_{(\zeta)} - \delta\mbf{x}_{(\xi)}\right),
  \label{eq:BD_identity}
\end{equation}
we obtain
\begin{align}
  \left.\delta\phi\right|_{\mbf{X}_s}
  &= \delta\phi_f + \delta\phi_g - \nabla_{\mbf{x}}\phi_f \cdot \left(\delta\mbf{x}_f - \delta\mbf{x}_s\right) - \nabla_{\mbf{x}}\phi_g \cdot \left(\delta\mbf{x}_g - \delta\mbf{x}_s\right) \notag \\
  &= \delta\phi_f + \delta\phi_g - \nabla_{\mbf{x}}\phi_f \cdot \delta\mbf{x}_f - \nabla_{\mbf{x}}\phi_g \cdot \delta\mbf{x}_g + \nabla_{\mbf{x}}\phi \cdot \delta\mbf{x}_s.
  \label{eq:delta_phi}
\end{align}
Defining
\begin{equation}
  P^\star_\phi \equiv \sum_{\xi=A,B} \rho_\xi\,\left.\p{\psi_\xi}{\phi}\right\vert_{\mbf{F}_s},
  \label{eq:Pstar_phi}
\end{equation}
and substituting~\eqref{eq:delta_phi} into~\eqref{eq:deltaU_solid}, the porosity dependent part of the solid energy variation becomes:
\begin{align}
  \int_{\mc{B}_{0}} \sum_{\xi=A,B} \rho_{\xi 0}\,\left.\p{\psi_\xi}{\phi}\right\vert_{\mbf{F}_s}\, \left.\delta\phi\right|_{\mbf{X}_s}\, dV_s
  &= \int_{\mc{B}} P^\star_\phi\, \delta\phi_f \,\dv + \int_{\mc{B}} P^\star_\phi\, \delta\phi_g \,\dv \notag \\
  &\quad - \int_{\mc{B}} P^\star_\phi\, \nabla_{\mbf{x}}\phi_f \cdot \delta\mbf{x}_f \,\dv - \int_{\mc{B}} P^\star_\phi\, \nabla_{\mbf{x}}\phi_g \cdot \delta\mbf{x}_g \,\dv \notag \\
  &\quad + \int_{\mc{B}} P^\star_\phi\, \nabla_{\mbf{x}}\phi \cdot \delta\mbf{x}_s \,\dv.
  \label{eq:phi_variation_split}
\end{align}

\subsection{Variation of the interfacial energy}

With $S_f = \phi_f/\phi$ and $\phi = \phi_f + \phi_g$, define the capillary pressure by
\begin{equation}
  p_c \equiv -\frac{1}{\phi}\frac{\partial\gamma_{fg}}{\partial S_f} > 0.
  \label{eq:pc_def}
\end{equation}
Then the variation of the interfacial energy contributes
\begin{align}
  \delta U_{fg}
  &= \int_{\mc{B}} \left[-p_c\frac{\phi_g}{\phi}\,\delta\phi_f + p_c\frac{\phi_f}{\phi}\,\delta\phi_g\right] \,\dv \notag \\
  &\quad + \int_{\mc{B}} \left[-p_c\frac{\phi_g}{\phi}\,\nabla_{\mbf{x}}\phi_f \cdot \delta\mbf{x}_f + p_c\frac{\phi_f}{\phi}\,\nabla_{\mbf{x}}\phi_g \cdot \delta\mbf{x}_g\right] \,\dv.
  \label{eq:deltaU_fg}
\end{align}

\subsection{Collected variational principle}

Evaluating all variations and collecting by independent variations $\delta\mbf{x}_\xi$, $\delta\phi_\xi$, $\delta\bar\rho_\xi$, and $\delta c_\xi$ gives
\begin{align}
  0 = \sum_\xi \int_{\mc{B}} \bigg\{
  &\bigg[-\rho_\xi\mbf{a}_\xi + \nabla_{\mbf{x}}\cdot(\bs\sigma'_\xi)^T + \rho_\xi\mbf{g} + \mbf{b}_\xi
  - \nabla_{\mbf{x}}\left(\mu_\xi J_\xi\right) + \lambda\nabla_{\mbf{x}}\phi_\xi + c^+_\xi\nabla_{\mbf{x}}\tau
  \bigg]\cdot\delta\mbf{x}_\xi \notag \\
  &+ \bigg[\frac{\mu_\xi J_\xi}{\phi_\xi} - \lambda\bigg] \,\delta\phi_\xi
  + \bigg[\frac{\mu_\xi J_\xi}{\bar\rho_\xi} - \rho_\xi\p{\psi_\xi}{\bar\rho_\xi}\bigg] \,\delta\bar\rho_\xi \notag \\
  &+ \bigg[\frac{1}{2}\mbf{v}_\xi\cdot\mbf{v}_\xi + \mc{L}_\xi - \frac{D_\xi\tau}{Dt} - \frac{\mu_\xi J_\xi}{\rho_\xi}\bigg]\frac{\delta c_\xi}{J_\xi}
  \bigg\} \,\dv \notag \\
  &+ \int_{\mc{B}} P^\star_\phi\, \delta\phi_f \,\dv + \int_{\mc{B}} P^\star_\phi\, \delta\phi_g \,\dv
  + \int_{\mc{B}} P^\star_\phi\, \nabla_{\mbf{x}}\phi \cdot \delta\mbf{x}_s \,\dv \notag \\
  &- \int_{\mc{B}} P^\star_\phi\, \nabla_{\mbf{x}}\phi_f \cdot \delta\mbf{x}_f \,\dv - \int_{\mc{B}} P^\star_\phi\, \nabla_{\mbf{x}}\phi_g \cdot \delta\mbf{x}_g \,\dv \notag \\
  &+ \int_{\mc{B}} \left[-p_c\frac{\phi_g}{\phi}\,\delta\phi_f + p_c\frac{\phi_f}{\phi}\,\delta\phi_g - p_c\frac{\phi_g}{\phi}\,\nabla_{\mbf{x}}\phi_f \cdot \delta\mbf{x}_f + p_c\frac{\phi_f}{\phi}\,\nabla_{\mbf{x}}\phi_g \cdot \delta\mbf{x}_g\right] \,\dv.
  \label{eq:collected}
\end{align}
This is the fully collected first variation.  The final two lines contain the new contributions from the porosity-dependent solid energy and the fluid--gas interfacial energy.

\subsection{Euler--Lagrange equations}

Since the variations are independent and arbitrary, each coefficient in~\eqref{eq:collected} must vanish.

\paragraph{Pressure identifications.}  From the coefficient of $\delta\bar\rho_\xi$ we identify the thermodynamic phase pressures
\begin{equation}
  \bar{p}_\xi \equiv \bar\rho_\xi^2\frac{\partial\psi_\xi}{\partial\bar\rho_\xi}, \qquad \xi \in \{f,g\},
  \label{eq:pbar_def}
\end{equation}
so that
\begin{equation}
  \frac{\mu_f J_f}{\phi_f} = \bar p_f, \qquad \frac{\mu_g J_g}{\phi_g} = \bar p_g.
  \label{eq:mu}
\end{equation}
From $\delta\phi_A$ and $\delta\phi_B$,
\begin{equation}
  \frac{\mu_A J_s}{\phi_A} = \lambda, \qquad \frac{\mu_B J_s}{\phi_B} = \lambda.
\end{equation}
From $\delta\phi_f$ and $\delta\phi_g$, including the $P^\star_\phi$ and capillary contributions,
\begin{align}
  \bar p_f &= \lambda - P^\star_\phi - p_c(1-S_f),
  \label{eq:pf}
  \\
  \bar p_g &= \lambda - P^\star_\phi + p_c S_f.
  \label{eq:pg}
\end{align}
Therefore
\begin{equation}
  \lambda = S_f\bar p_f + (1-S_f)\bar p_g + P^\star_\phi.
  \label{eq:lambda}
\end{equation}

\paragraph{Fluid momentum.}  From the coefficient of $\delta\mbf{x}_f$,
\begin{equation}
  \rho_f\mbf{a}_f = \rho_f\mbf{g} + \mbf{b}_f - \nabla_{\mbf{x}}(\mu_f J_f) + \lambda\nabla_{\mbf{x}}\phi_f - P^\star_\phi\nabla_{\mbf{x}}\phi_f - p_c\frac{\phi_g}{\phi}\nabla_{\mbf{x}}\phi_f + c^+_f\nabla_{\mbf{x}}\tau.
  \label{eq:mom_fluid_raw}
\end{equation}
Using $\mu_f J_f = \phi_f \bar p_f$ and \eqref{eq:pf},
\begin{align*}
  -\nabla_{\mbf{x}}(\phi_f \bar p_f) + \lambda\nabla_{\mbf{x}}\phi_f - P^\star_\phi\nabla_{\mbf{x}}\phi_f - p_c\frac{\phi_g}{\phi}\nabla_{\mbf{x}}\phi_f
  &= -\phi_f\nabla_{\mbf{x}}\bar p_f - \bar p_f\nabla_{\mbf{x}}\phi_f \\
  &\quad + \left(\lambda - P^\star_\phi - p_c\frac{\phi_g}{\phi}\right)\nabla_{\mbf{x}}\phi_f \\
  &= -\phi_f\nabla_{\mbf{x}}\bar p_f.
\end{align*}
Hence the fluid momentum balance reduces to the phase-pressure form
\begin{equation}
  \boxed{\rho_f\mbf{a}_f = \rho_f\mbf{g} + \mbf{b}_f - \phi_f\nabla_{\mbf{x}}\bar p_f + c^+_f\left(\nabla_{\mbf{x}}\tau - \mbf{v}_f\right),}
  \label{eq:mom_fluid}
\end{equation}
where the reaction-momentum term has been written in the standard Bedford--Drumheller form.

\paragraph{Gas momentum.}  From the coefficient of $\delta\mbf{x}_g$,
\begin{equation}
  \rho_g\mbf{a}_g = \rho_g\mbf{g} + \mbf{b}_g - \nabla_{\mbf{x}}(\mu_g J_g) + \lambda\nabla_{\mbf{x}}\phi_g - P^\star_\phi\nabla_{\mbf{x}}\phi_g + p_c\frac{\phi_f}{\phi}\nabla_{\mbf{x}}\phi_g + c^+_g\nabla_{\mbf{x}}\tau.
  \label{eq:mom_gas_raw}
\end{equation}
Using $\mu_g J_g = \phi_g \bar p_g$ and \eqref{eq:pg},
\begin{align*}
  -\nabla_{\mbf{x}}(\phi_g \bar p_g) + \lambda\nabla_{\mbf{x}}\phi_g - P^\star_\phi\nabla_{\mbf{x}}\phi_g + p_c\frac{\phi_f}{\phi}\nabla_{\mbf{x}}\phi_g
  &= -\phi_g\nabla_{\mbf{x}}\bar p_g - \bar p_g\nabla_{\mbf{x}}\phi_g \\
  &\quad + \left(\lambda - P^\star_\phi + p_c\frac{\phi_f}{\phi}\right)\nabla_{\mbf{x}}\phi_g \\
  &= -\phi_g\nabla_{\mbf{x}}\bar p_g.
\end{align*}
Hence the gas momentum balance reduces to
\begin{equation}
  \boxed{\rho_g\mbf{a}_g = \rho_g\mbf{g} + \mbf{b}_g - \phi_g\nabla_{\mbf{x}}\bar p_g + c^+_g\left(\nabla_{\mbf{x}}\tau - \mbf{v}_g\right).}
  \label{eq:mom_gas}
\end{equation}

\paragraph{Overall momentum.}  Summing the solid, fluid, and gas balances gives the mixture equation.  The direct inter-component body-force pair terms from~\eqref{eq:phi_variation_split} cancel in the sum, but the porosity-conjugate contribution survives through the solid pressure algebra.  Writing the weighted phase pressure as $\bar p \equiv S_f\bar p_f + (1-S_f)\bar p_g$, the overall balance takes the form
\begin{equation}
  \boxed{\rho\mbf{a} = \nabla_{\mbf{x}}\cdot(\bs\sigma'_s)^T - \nabla_{\mbf{x}}\bar p - \nabla_{\mbf{x}}\left[(1-\phi)P^\star_\phi\right] + \rho\mbf{g} - c^+_f\left(\mbf{v}_f - \mbf{v}_s\right) - c^+_g\left(\mbf{v}_g - \mbf{v}_s\right).}
  \label{eq:mom_overall}
\end{equation}
When $\psi_\xi$ does not depend on $\phi$ ($P^\star_\phi = 0$), this reduces to the standard mixture form.

\paragraph{Evolution of $\tau$.}  From the coefficient of $\delta c_\xi$ and choosing to evaluate the functions for the reactant $A$,
\begin{equation}
  \boxed{\left.\p{\tau}{t}\right|_{\mbf{X}_s} = \frac{1}{2}\mbf{v}_s\cdot\mbf{v}_s - \mc{L}_A - \frac{\lambda}{\bar\rho_A}.}
  \label{eq:tau_final}
\end{equation}

\subsection{Connection to the Coleman--Noll procedure}
\label{sec:coleman_connection}

The variational derivation and the Coleman--Noll procedure are complementary.  Hamilton's principle takes the free energies (e.g. $\psi_\xi(\mbf{F}_s, \phi)$) as functions of \emph{independent fields} (subject to constraints) and derives the momentum balances, pressure identifications, and body forces as Euler--Lagrange equations.  The Coleman--Noll procedure (Appendix~\ref{sec:coleman}) starts from the entropy inequality with the \emph{same} energy and derives constitutive restrictions.  The multiplicative decomposition $\mbf{F}_s = \mbf{A}_s \bar{\mbf{F}}_s$ is a constitutive choice imposed within the Coleman--Noll procedure, giving $\psi_\xi(\mbf{F}_s, \phi) = \psi^{\text{iso}}_\xi(\bar{\mbf{F}}_s) + \psi^{\text{dis}}_\xi(J_s, \phi)$, which is a specialization of the general $\psi_\xi(\mbf{F}_s, \phi)$ used in the variational principle.
